% Apêndice A — Reprodutibilidade (orçamento: ~6 pp)
% A matriz de evidência é a peça mais importante: lista as afirmações, incluindo as retiradas.

\chapter{Reprodutibilidade}
\label{ap:reprodutibilidade}

\section{Ambiente}
\label{sec:ap_ambiente}

O sistema está escrito em Python~3.12. As dependências do ambiente que produz os números desta
dissertação estão fixadas num ficheiro de bloqueio, para que esse ambiente possa ser recriado noutra
máquina. O serviço implantado tem o seu próprio conjunto, mais pequeno, porque não precisa de treinar
nada. As versões com que todos os números deste documento foram produzidos são
\texttt{numpy}~2.1.3, \texttt{pandas}~2.2.3, \texttt{scikit-learn}~1.9.0,
\texttt{matplotlib}~3.11.0, \texttt{yfinance}~1.4.1, \texttt{sentence-transformers}~5.6.0,
\texttt{transformers}~5.12.1 e \texttt{torch}~2.12.1 na variante para CPU. Estão fixadas, e não em
intervalos, porque uma biblioteca que muda sozinha muda os resultados sem aviso.


Não é preciso GPU. Tudo o que esta dissertação reporta corre num computador portátil comum, o que é
uma consequência deliberada das escolhas do Capítulo~\ref{cap:metodologia}.

A verificação automática do sistema tem \textbf{763 testes}, entre os quais os que impõem a
separação temporal descrita na Secção~\ref{sec:met_avaliacao}.

\section{Como cada número é reproduzido}
\label{sec:ap_numeros}

Nenhum \emph{resultado} desta dissertação foi escrito à mão. Cada um é produzido por um
procedimento automático com semente fixa, que escreve o seu próprio ficheiro de resultados; o texto
cita esse ficheiro. A Tabela~\ref{tab:ap_numeros} reúne todos os resultados reportados e indica onde
cada um aparece no corpo do documento; o ficheiro de resultados correspondente, e o procedimento que
o escreve, estão nomeados no repositório de código.

A garantia de determinismo vale para tudo o que está acima da linha divisória da tabela. As três
linhas de baixo são de outra natureza e por isso estão separadas: são medições sobre a
\emph{operação}, e o seu valor depende de quando o registo foi lido, não de uma semente. O capítulo
onde aparecem diz sempre sobre que período e que conjunto foram feitas.

Os restantes números do documento são de três tipos. Os \textbf{valores retirados das fontes
citadas} estão citados no sítio onde aparecem. As \textbf{constantes de configuração} estão na
Tabela~\ref{tab:sis_portas}, com a distinção entre as que foram derivadas e as que foram
escolhidas. E há um terceiro tipo, que convém não esconder: \textbf{medições minhas que ficam
fora desta tabela} porque servem uma passagem e não um resultado, como as $820$ linhas que o
embargo descarta, o $\sigma$ do caso trabalhado da Tesla, o custo de memória dos três formatos, ou
as contagens por fonte de notícias. Cada uma diz, no sítio onde aparece, sobre que dados foi feita.

Uma delas merece nota própria, porque é a única que não se pode regenerar. A repartição do funil
por porta, na Tabela~\ref{tab:sis_funil}, foi lida do registo de decisões num dia concreto, e esse
registo guarda apenas três dias: a razão está na Secção~\ref{sec:met_dados}, e é que ele é
republicado a cada ciclo, portanto o custo é de publicação e não de armazenamento. Correr o mesmo
procedimento hoje daria outro dia, não o mesmo. Fica dito assim, em vez de se apresentar como
reproduzível uma medição que o não é.

Declarar que um número não se regenera não é, porém, o mesmo que mostrar de onde ele veio, e essa
segunda parte faltava: a repartição estava escrita apenas nesta dissertação. Passou a estar
também num ficheiro de evidência, ao lado dos outros, com a data da leitura e com o procedimento
que a produziu, de forma a que qualquer dia futuro tenha origem mesmo que aquele já não volte.

\begin{table}[H]
\caption[Cada resultado e o procedimento que o produz]{Todos os resultados reportados, e onde são
gerados. Correr outra vez sobre os mesmos dados devolve exatamente os mesmos valores.}
\label{tab:ap_numeros}
\centering
\small
\begin{tabular}{p{0.46\textwidth} L{0.20\textwidth} L{0.22\textwidth}}
\toprule
\tabhead{Resultado} & \tabhead{Valor} & \tabhead{Onde aparece} \\
\midrule
Amplitude da taxa de disparo ($z$-score vs limiar fixo) & $0.015$ vs $0.344$ & §\ref{sec:av_qi1} \\
$F_1$ contra o rótulo aproximado & $0.516$ vs $0.218$ & §\ref{sec:av_qi1} \\
$z$-score contra dois detetores aprendidos & $0.530$ / $0.269$ / $0.280$ & §\ref{sec:av_qi1} \\
Precisão@5 da recuperação (4 métodos) & $0.514$ / $0.346$ / $0.240$ / $0.126$ & §\ref{sec:av_qi2} \\
Precisão@5 no protocolo simétrico de escala & $0.595$ com chão $0.333$ & §\ref{sec:av_qi2} \\
Precisão@5 só com precedentes anteriores & $0.513$ com chão $0.259$ &
§\ref{sec:av_qi2_causal} \\
Concordância de direção vs chão de acaso & $0.708$ vs $0.688$ & §\ref{sec:av_qi2} \\
Triagem, PR-AUC das seis famílias & $0.378$ a $0.542$ & §\ref{sec:av_qi3} \\
Precisão dentro do orçamento (4 formas de escolher) & $0.163$ / $0.379$ / $0.662$ / $0.632$ &
§\ref{sec:av_qi3} \\
Grelha de nove definições de rótulo & $9$ de $9$ & §\ref{sec:av_qi3} \\
Decomposição de um movimento, caso trabalhado & $-0.3992$ / $-0.0362$ / $+6.7298$ \% &
§\ref{sec:met_decomposicao} \\
$R^2$ da decomposição, mediana e casos negativos & $0.460$, $1$ de $17$ &
§\ref{sec:av_decomposicao} \\
Ablação da identidade: tabela de consulta vs modelo & $0.534$ vs $0.538$ &
§\ref{sec:av_ablacao} \\
O texto por cima da tabela de consulta & $0.547$ contra $0.534$ &
§\ref{sec:av_acrescimo} \\
Selectividade da porta, dentro vs entre empresas & $0.064$ vs $0.385$ & §\ref{sec:av_portao} \\
Deriva de distribuição, \gls{PSI} da volatilidade & $0.281$ & §\ref{sec:av_deriva} \\
Sistema inteiro contra as alternativas & $0.489$ / $0.632$ / $0.662$ / $0.968$ &
§\ref{sec:av_pontaaponta} \\
Ablação à janela e ao peso do desvio & $0.678$ e $0.664$ contra $0.516$ &
§\ref{sec:av_alternativas} \\
\midrule
\multicolumn{3}{l}{\emph{Medições sobre a operação, sem semente: dependem de quando o registo foi
lido}} \\
Decisões maturadas em produção, mantidas vs suprimidas & $0.589$ vs $0.617$ &
§\ref{sec:av_qi3} \\
Cobertura das notícias em dias invulgares & $88.5\%$ & §\ref{sec:con_limitacoes} \\
Latência, por era e decomposta & $196$ / $143$ / $158$ min & §\ref{sec:sis_implantacao} \\
Funil de produção (captadas $\rightarrow$ enviadas) & $944 \rightarrow 42$ & §\ref{sec:sis_portas} \\
Funil de um dia, repartido por porta & $5\,060 \rightarrow 5$ & §\ref{sec:sis_portas} \\
\bottomrule
\end{tabular}
\end{table}

\section{Matriz de evidência}
\label{sec:ap_matriz}

A tabela anterior lista \emph{resultados}. Esta lista \emph{afirmações}, que é um teste mais duro,
porque uma afirmação pode sobreviver à evidência que a sustentava.

A Tabela~\ref{tab:ap_matriz} põe cada afirmação substantiva desta dissertação frente à evidência que
a suporta e ao estado em que ficou. \textbf{As afirmações retiradas ficam na tabela em vez de serem
apagadas}, porque uma matriz que só lista as sobreviventes não é uma auditoria.

\begin{table}[H]
\caption[Matriz de evidência]{Cada afirmação, a evidência, e o que lhe aconteceu. ``Retirada''
significa que este trabalho a abandonou depois de a medir.}
\label{tab:ap_matriz}
\centering
\scriptsize
\renewcommand{\arraystretch}{1.15}
\begin{tabular}{p{0.36\textwidth} p{0.36\textwidth} L{0.16\textwidth}}
\toprule
\tabhead{Afirmação} & \tabhead{Evidência} & \tabhead{Estado} \\
\midrule
\multicolumn{3}{l}{\textbf{QI1, deteção}} \\
Dispara a um ritmo comparável entre empresas muito diferentes & Amplitude $0.015$ vs $0.344$, sem
usar rótulos & Mantida \\
Bate um limiar fixo no rótulo aproximado & $F_1$ $0.516$ vs $0.218$ & Mantida (o rótulo é um proxy) \\
Bate dois detetores aprendidos com a mesma informação & $F_1$ $0.530$ vs $0.269$ e $0.280$ & Mantida \\
Não usa informação do futuro & A janela exclui o dia julgado; imposto por teste & Mantida \\
\midrule
\multicolumn{3}{l}{\textbf{QI2, precedentes}} \\
Bate as alternativas lexical e triviais & P@5 $0.514$ vs $0.346$ / $0.240$ / $0.126$ & Mantida \\
Bate por larga margem o melhor chão trivial & ``Sempre o setor maior'' dá $0.467$: a margem é
$+0.047$ e não $+0.274$ & \textbf{Estreitada} (o agregado media a composição do corpus; a
afirmação passa a ser por setor, onde ela se mantém nos cinco) \\
Aguenta no protocolo simétrico de escala & P@5 $0.595$ com chão $0.333$; permite candidatos
posteriores & Mantida como teste de escala, não como tarefa causal \\
Aguenta a restrição da produção (só o passado) & P@5 $0.513$ com chão $0.259$: a margem sobre
o acaso muda $0.008$ & \textbf{Estreitada} (o número simétrico descrevia uma tarefa mais fácil) \\
Capta tema, não direção & Concordância $0.708$ contra um chão de $0.688$ & Mantida (é uma limitação) \\
``Relevante'' é o mesmo que ``mesmo setor'' & --- & \textbf{Não afirmado} (é um proxy) \\
O chão de semelhança escolhe precedentes que preveem melhor & Acima do chão $0.504$, abaixo $0.506$,
acaso $0.507$ & \textbf{Retirada} \\
\midrule
\multicolumn{3}{l}{\textbf{QI3, triagem}} \\
Nenhum modelo com texto bate a volatilidade & PR-AUC $0.496$ vs $0.542$; aguenta re-teste justo,
reamostragem por grupos e nove rótulos & Mantida \\
O texto acrescenta por cima da tabela de consulta & PR-AUC $0.547$ contra $0.534$: $+0.012$,
IC $[+0.004, +0.020]$ por grupos & \textbf{Nova} (positiva, e não muda a precisão no orçamento) \\
Sem uso de informação futura nas entradas & Mutar o futuro deixa as entradas iguais e muda o rótulo &
Mantida \\
A triagem quase quadruplica a precisão no orçamento & O chão de $0.163$ ordena alfabeticamente; ao
acaso dá $0.379$ & \textbf{Retirada} (o ganho é $1.67\times$) \\
É a aprendizagem que produz o ganho no orçamento & Um prior de $13$ constantes dá $0.662$ contra
$0.632$ do modelo & \textbf{Retirada} \\
O modelo implantado seleciona melhor do que o acaso & $0.589$ mantidas contra $0.617$ suprimidas &
\textbf{Retirada} \\
\midrule
\multicolumn{3}{l}{\textbf{Sistema}} \\
O texto do alerta não pode divergir do cálculo & É montado a partir dos objetos calculados & Mantida \\
A fonte gratuita limita a cobertura & Pelo menos uma notícia em $88.5\%$ dos dias invulgares &
Mantida (limite superior) \\
As explicações são \emph{úteis} a uma pessoa & --- & \textbf{Não afirmado} \\
\bottomrule
\end{tabular}
\end{table}

Duas observações sobre a tabela no seu conjunto, e não sobre nenhuma linha.

Primeiro: \textbf{quatro afirmações foram retiradas, e todas o foram pelas medições deste próprio
trabalho}, nenhuma por revisão externa. É o comportamento pretendido de uma avaliação desenhada
antes das experiências: tem de ser capaz de dizer que não.

Segundo: as duas linhas marcadas ``não afirmado'' são tão importantes como as outras. Uma delas é a
utilidade para o utilizador, que não foi medida porque não houve estudo com pessoas. Onde não há
evidência, não há afirmação.

\section{Dados e atribuição}
\label{sec:ap_dados}

Os corpora de origem não são versionados: o conjunto histórico completo ocupa centenas de megabytes
e é regenerado pelos procedimentos de preparação. Acompanham o código pequenas amostras, para que os
testes corram sem rede, e uma excepção que convém declarar: a \textbf{base de casos que o sistema
consulta} está versionada, com $56$ megabytes de vetores e $8$ de metadados. Não é uma amostra, é o
artefacto de que a aplicação implantada precisa para arrancar, e sem ele o sistema não teria
precedentes para mostrar.

O conjunto histórico usado para construir a base de casos é o \gls{FNSPID}
\autocite{dong2024fnspid}, utilizado ao abrigo da sua licença Creative Commons BY-SA 4.0, cuja
atribuição é aqui cumprida. Apenas excertos mínimos de títulos aparecem no documento, a título
ilustrativo.

Essa licença tem uma segunda obrigação que convém não calar, porque é acionada: como o repositório
distribui ficheiros \emph{derivados} do \gls{FNSPID}, a cláusula de partilha nos mesmos termos
aplica-se a eles. A licença do código é uma decisão que fica em aberto neste documento, e é esta a
restrição que a condiciona.

Os preços vêm de serviços gratuitos de dados de mercado, através de uma cadeia de alternativas
descrita na Secção~\ref{sec:sis_dados}; as notícias em tempo real vêm de uma API gratuita, dentro dos
seus limites de utilização. Sobre dados pessoais convém ser exato em vez de generoso. O sistema
\textbf{não sabe o que os utilizadores possuem e não lhes pergunta}, que é a decisão de desenho que
interessa: não há carteira, não há posições, não há montantes. O canal público não identifica
ninguém. Há, no entanto, uma funcionalidade opcional em que quem quiser receber alertas
personalizados fala com o robô, e aí ficam guardados o identificador de conversa do Telegram e a
lista de empresas que essa pessoa escolheu seguir. É pouco, e é a totalidade: fica dito, porque
``nenhum dado pessoal'' seria mais simples de escrever e não seria verdade.

\section{O estudo com utilizadores: desenho}
\label{sec:ap_estudo}

A limitação mais séria deste trabalho é não ter sido testado com utilizadores. O
Capítulo~\ref{cap:conclusoes} refere um protocolo já preparado, e o seu desenho fica aqui para que
essa afirmação possa ser conferida.

O estudo compara duas condições sobre seis alertas reais retirados do canal: a condição~A apresenta
o facto sem explicação, a condição~B apresenta o alerta completo. Dois dos seis são casos em que o
tema e a direção do movimento discordam, que é onde a Secção~\ref{sec:av_qi2} mostra o sistema ser
mais fácil de interpretar mal. O contrabalanço cruza dois fatores: a ordem das condições e qual das
metades dos estímulos serve de condição~A. O segundo é necessário porque as condições têm de usar
alertas diferentes, sob pena de a segunda medir memória; sem o cruzamento, qualquer diferença entre
condições ficaria confundida com uma das metades ser mais fácil.

A pergunta é a que a verificação automática não alcança. A fidelidade entre cada frase do alerta e
o facto que a sustenta está garantida por construção, mas a afirmação do produto é uma travessia
que um leitor faz: dada uma frase com âncora, consegue abrir o facto e julgar se ele a sustenta,
sozinho? Se não conseguir, a contribuição é verdadeira e inútil. É uma pergunta qualitativa e não
exige amostra grande.

Duas salvaguardas metodológicas estão fixadas no código antes de existirem dados: o limiar de oito
participantes para o teste estatístico, cuja alteração ficaria registada no histórico, e o
procedimento de análise, que sobre a folha por preencher não produz resultado nenhum. Um terceiro
bloco do desenho original, sobre texto gerado, está marcado como não executável, porque as rotas
que o serviam foram retiradas do produto.
